\chapter{A Verse on Layout}	\OnePageChapter         % ONE PAGE!

\noindent\paragraph{(Written for this example)}
(A one-page chapter --- page must be numbered!)
This appendix shows how a short, verse-like block sets
inside a one-page chapter. The lines below were written
for this sample document and quote no other work; their
subject is the layout that the class itself provides.


\begin{quotation}
\noindent The page begins with level margins, quiet, wide and white, \\
a frame that holds the argument and keeps the columns right. \\
The chapters open, numbered, on a fresh and empty sheet; \\
the sections follow, step by step, in order clean and neat. \\
Each figure waits beside its words, its caption placed below; \\
each table wears its caption high, the first thing it will show. \\
The footnotes speak in smaller type beneath the running prose, \\
and citations send the reader back to where each claim arose. \\
The contents list the road ahead, the pages mark the way, \\
and appendices hold the rest the chapters could not say. \\
When every line is in its place and every number true, \\
the document is finished, and it waits for someone new.
\end{quotation}


\noindent(A short prose note between the two blocks shows how
the class returns to body text after a quotation and then
sets a second one.)


\begin{quotation}
\noindent So here's to every rule that keeps the layout in its place: \\
the spacing and the numbering, the headings and the face. \\
You never read a single line, nor weigh what has been planned, \\
yet every thesis leans on you, a steady, silent hand.
\end{quotation}
